\def\stat{garanin}





\def\tit{МЕТОДИЧЕСКИЕ ВОПРОСЫ ФОРМИРОВАНИЯ СИСТЕМЫ ТЕХНИЧЕСКОГО ОБЕСПЕЧЕНИЯ 


ИНФОРМАЦИОННО-ТЕЛЕКОММУНИКАЦИОННЫХ СЕТЕЙ}





\def\titkol{Методические вопросы формирования системы технического обеспечения 


ИТС} %информационно-телекоммуникационных сетей} 





\def\autkol{А.\,А.~Зацаринный, Н.\,Г.~Буроменский, А.\,И.~Гаранин}





\def\aut{А.\,А.~Зацаринный$^1$, Н.\,Г.~Буроменский$^2$, А.\,И.~Гаранин$^3$}





\titel{\tit}{\aut}{\autkol}{\titkol}





\renewcommand{\thefootnote}{\arabic{footnote}}


\footnotetext[1]{Институт проблем информатики Российской академии наук, 


azatsarinny@ipiran.ru} 


\footnotetext[2]{Филиал ФГКУ <<46 ЦНИИ>> Минобороны России, г.\ Мытищи, buromienskii@mail.ru} 


\footnotetext[3]{Институт проблем информатики Российской академии наук, algaranin@mail.ru} 





            \label{st\stat}


            


            \vspace*{-12pt}





  \Abst{Приведены основные понятия применительно к объекту исследования. Показано, 


что формирование системы технического обеспечения (СТО) 


  информационно-телекоммуникационных сетей (ИТС) является комплексной 


системотехнической проблемой, охватывающей весьма широкий круг взаимосвязанных 


процессов на всех этапах жизненного цикла ИТС. Сформулированы основные из них 


примерно в той последовательности, в которой они должны решаться в ходе создания СТО. 


Последовательно излагаются этапы проведения исследований по созданию СТО, задачи и 


получаемые при этом результаты. Дана аналитическая модель функционирования СТО и 


общая схема метода формирования СТО. Показано, что математическое описание процессов 


функционирования сложных иерархических систем, подобных системе СТО ИТС,~--- 


чрезвычайно сложная и серьезная проблема. Решать ее предлагается в два этапа. На первом 


этапе описание дается в форме абстрактной математической формализации процесса 


функционирования с использованием теории множеств (отношений на множестве или между 


множествами и функциональных отображений). На втором этапе выделяются отдельные 


наиболее важные процессы системы, такие как ремонт, техническое обслуживание, 


снабжение новыми техническими средствами (ТС), ЗИП (запасные части, инструменты, принадлежности), 


и осуществляется переход от 


интуитивных понятий к детализированному математическому описанию с использованием 


аналитико-вероятностных моделей. Предложена примерная структурная схема СТО. 


Обоснован базовый состав нормативного обеспечения СТО ИТС. }


  


  \KW{система; информационно-телекоммуникационная система; система технического 


обеспечения; исследование; проектирование; разработка; эксплуатация; эффективность; 


моделирование; структура}





\DOI{10.14357\!/08696527130211}


  





%  \vskip 14pt plus 9pt minus 6pt








      \thispagestyle{headings}





      





            


  %          \vspace*{-12pt}





  


  \section{Введение}


  В настоящее время идет процесс широкомасштабного насыщения ИТС 


современными сложными ТС, к которым относятся 


прежде всего средства передачи, приема, обработки и отображения 


информации. Вместе с тем известно, что чем сложнее ТС, 


тем оно менее надежно и, как следствие, подвержено частым отказам. Однако 


причиной низкой на\-деж\-ности ИТС и ее элементов часто является не только их 


сложность, но и неполное обеспечение условий эксплуатации (УЭ) ИТС, низкая 


автоматизация процессов управления системой технического обеспечения ИТС, 


слабая подготовка обслуживающего персонала. 


  


  В связи с этим одним из направлений поддержания ИТС в постоянной 


готовности является создание эффективной СТО. Под СТО ИТС будем 


понимать комплекс мероприятий по укомплектованию элементов ИТС 


техническими средствами, поддержание их в постоянной готовности к 


применению по назначению и своевременному восстановлению отказавших 


элементов. 


  


  Необходимо отметить, что работы по проблеме разработки СТО практически 


отсутствуют. Причиной этого положения являются такие особенности СТО, как 


ее малосерийность и уникальность. И~в то же время значительное число работ 


посвящено методам оценки надежности сетей связи~[1--4] и управлению 


эффективностью процесса технической эксплуатации отдельных видов 


техники~[5,~6]. 


  


  Данная работа посвящена разработке методологических основ формирования 


СТО ИТС и является дальнейшим развитием идей и результатов, изложенных 


в~[7]. Дана общая последовательность работ при проведении исследований, 


проектировании, разработке и эксплуатации. Дана аналитическая и общая 


модель формирования СТО ИТС. При разработке методологических основ 


формирования СТО ИТС использованы основные положения, термины и 


определения, приведенные в нормативных документах~[8--11].


  


\section{Основное содержание исследований при~разработке системы~технического обеспечения


информационно-телекоммуникационных~сетей}


  


  Согласно~\cite{8-gr}, весь период существования образца техники именуется 


процессом его жизненного цикла. Первым процессом жизненного цикла 


являются исследования. Исследования по созданию СТО с целью большей 


конкретизации целевых задач разделим на три этапа. 


  


  В ходе \textbf{первого этапа} исследования должны быть решены следующие 


во\-просы:


  \begin{itemize}


  \item анализ проблем создания СТО;


  \item описание объекта исследования;


  \item определение структуры СТО;


  \item разработка требований к системе;


  \item формулировка критерия эффективности функционирования системы.


  \end{itemize}


  


  Целью \textbf{второго этапа} является анализ условий применения СТО. 


  


  Содержание задач, решаемых на этом этапе, состоит в том, что вся 


совокупность элементов и процессов, имеющих отношение к поставленной 


цели, разбивается на следующие группы:


  \begin{itemize}


  \item собственно СТО;


  \item надсистема~--- информационно-телекоммуникационная система;


  \item внешняя среда.


  \end{itemize}


  


  Результатом этого этапа является определение состава и уровней факторов, 


воздействующих на элементы ИТС, элементы системы СТО и оценка 


последствий такого воздействия. 


  


  Формально УЭ СТО можно определить как <<все, что не СТО>>, но прямо 


или косвенно оказывает прямое воздействие на выполнение ею своих задач.


  


  Из числа параметров СТО, получаемых в ходе выполнения этого этапа 


исследований, можно выделить сравнительно небольшую, но очень важную 


группу параметров, таких как прогнозируемые потоки отказов элементов ИТС, 


распределение элементов ИТС по видам отказов, требования к среднему 


времени восстановления элементов ИТС. Эти характеристики будем называть 


внешними характеристиками, а этап их выбора и обоснования в литературе 


именуется внешним проектированием~\cite{12-gr}.


  


  Последовательность работ при внешнем проектировании, как правило, такая: 


  \begin{enumerate}[(1)]


  \item разработка моделей УЭ СТО;\\[-9pt]


  \item выбор критерия эффективности СТО;\\[-9pt]


  \item анализ зависимости критерия эффективности от УЭ;\\[-9pt]


  \item оценка потоков отказов элементов СТО;\\[-9pt]


  \item разработка методов ремонта;\\[-9pt]


  \item разработка требований к ТС восстановления 


элементов ИТС;\\[-9pt]


  \item выбор рациональных вариантов организации СТО.


  \end{enumerate}


  


  Следует отметить, что эффективность СТО ИТС во многом зависит от 


применяемых видов ремонта и ТС восстановления 


отказавших элементов. При этом в ходе их разработки целесообразно 


использовать опыт создания СТО военной связи и автоматизированных 


средств управления~\cite{13-gr}, а также ремонта радиоэлектронного 


вооружения США~\cite{14-gr}.


  


  Информационную основу внешнего проектирования СТО 


применительно к УЭ составляют данные, включающие: 


  \begin{itemize}


\item прогнозируемые УЭ;


\item организационно-штатную структуру элементов ИТС; 


\item требования к укомплектованности элементов ИТС;


\item классификацию ТС по видам отказов~\cite{15-gr}.


\end{itemize}





  На \textbf{третьем этапе} производится разработка методов моделирования СТО. 


  


  Отметим, что математическое описание процессов функционирования 


иерархических систем, подобных системе СТО ИТС,~--- чрезвычайно сложная 


и серьезная проблема. В~настоящее время пока еще не существует единого 


математического аппарата, который бы позволил решить проблему описания 


функционирования сложных систем в полной мере, с достаточной строгостью и 


обоснованностью. 


  


  Поэтому решение задачи моделирования СТО ИТС выполним в следующей 


последовательности.





Сначала дадим описание процесса функционирования СТО ИТС в виде 


абстрактной математической модели с использованием теории множеств 


(отношений на множестве или между множествами, функциональных 


отображений и~т.\,д.).


  


  Далее выделим отдельные процессы системы (например, ремонт, 


техническое обслуживание, снабжение ЗИП и~т.\,д.)\ и осуществим переход от 


интуитивных понятий к детализированному математическому описанию с 


использованием общих аналитических~\cite{16-gr}, вероятностных и 


физических моделей. Последние применяют для выявления конкретных 


технических причин, вызывающих отказы ТС. Методические основы 


разработки физической модели при радиационном воздействии на 


радиоэлектронные системы изложены в работе~\cite{17-gr}, а особенности 


восстановления техники связи при ее радиационном заражении показаны на 


примере ликвидации аварии на Чернобыльской АЭС~\cite{18-gr}.





\section{Основное содержание работ при~проектировании и~разработке 


системы технического обеспечения}


  


  Общим результатом исследований должно стать так\-ти\-ко-тех\-ни\-че\-ское 


задание на проектирование и разработку СТО ИТС, а также отдельных ее 


элементов. Учитывая сложность создания СТО при ее проектировании и 


разработке в рамках опыт\-но-кон\-ст\-рук\-тор\-ских работ (ОКР), как правило, 


определяется головная организация и исполнители составных частей. При этом 


головная организация и исполнители составных частей должны планировать и 


управлять проектированием и разработкой системы в соответствии с 


требованиями нормативных документов~\cite{19-gr, 20-gr}.


  


  Анализ задач проектирования и разработки таких сложных изделий, как 


СТО, необходимо проводить с использованием методов системного анализа. 


В~соответствии с общей методологией исследований по разработке сложных 


систем можно выделить три основные группы факторов, определяющих выбор 


проектных решений:


  \begin{enumerate}[(1)]


  \item условия эксплуатации (взаимодействие с внешней средой);


  \item внутреннее строение (структура); 


  \item общесистемные, интегральные свойства.


  \end{enumerate}


  


  Важнейшими этапами проектирования и разработки являются верификация и 


валидация проекта системы. Верификация предполагает разработку и 


осуществление контроля или другой деятельности, необходимой для 


подтверждения соответствия разработанной системы СТО и (или) ее элементов 


требованиям заказчика. При этом деятельность по подтверждению требований, 


заданных в техническом задании (ТЗ) на ОКР, может подразумевать:


  \begin{itemize}


  \item  осуществление альтернативных расчетов;


  \item  сравнение спецификации на новый проект с аналогичной 


документацией на апробированный проект;


  \item проведение испытаний и демонстраций;


  \item анализ документов до их выпуска.


  \end{itemize}


  


  Валидация определяет подтверждение посредством объективных 


свидетельств того, что требования, предназначенные для конкретных УЭ СТО, 


выполнены.


  


  С учетом сложности разрабатываемой системы и элементов УЭ при 


валидации могут быть реальными или смоделированными. При этом по мере 


увеличения сложности растет доля смоделированных как элементов, так и УЭ. 


Моделируются прежде всего отказы элементов ИТС, восстановление 


работоспособности, измерение характеристик ТС, снабжение 


ЗИП.





Важно отметить и следующие методические аспекты создания СТО ИТС: 


\begin{enumerate}[1.]


  \item  При проектировании, изготовлении ТС, ремонте и комплектовании 


ЗИП должны быть установлены правила применения изделий электронной 


техники и электротехники иностранного производства в ТС ИТС и ТС СТО, а 


также закупки и оценки пригодности ТС. В качестве исходных при разработке 


таких правил можно рекомендовать работы~[21--23]. 


  \item  Система технического обеспечения является частью ИТС. Поэтому в ТЗ на создание ИТС заказчиком 


одновременно задаются требования и по техническому обеспечению. 


  Система технического обеспечения должна разрабатываться как составная часть ИТС и на всех этапах 


разработки (эскизное, техническое, рабочее проектирование) подвергаться всем 


видам испытаний совместно с элементами ИТС как на стенде главного 


конструктора, так и на объектах заказчика.


\end{enumerate}





\section{Методические рекомендации по~построению системы технического обеспечения 


информационно-телекоммуникационных сетей}


  


  Система СТО, как и любая сложная система, должна быть наделена вполне 


определенной структурой, под которой будем понимать совокупность 


элементов (подсистем) и множество устойчивых связей между ними. 


Особенностью формирования структуры СТО является учет территориальной 


распределенности ИТС. С этой точки зрения структуру СТО целесообразно 


разделить на относительно самостоятельные в функциональном отношении 


части (подсистемы). Подсистемы могут иметь региональный характер, однако 


строиться на единых для СТО принципах функционирования и управления. 


Рекомендуемая типовая структура СТО приведена на рис.~1. 


  


  Дадим небольшие пояснения структуры СТО. Выделение в ее структуре 


системы О$_0$ высшего уровня обусловлено тем, что используемые силы и 


средства являются лишь частью сил и средств, участвующих в организации 


функционирования ИТС. При этом система О$_0$ выполняет часть задач по 


техническому обеспечению в интересах нижнего звена управления.


     


     \begin{figure}[b] %fig1


     \vspace*{6pt}


\begin{center}


 \mbox{%


 \epsfxsize=115.8mm


 \epsfbox{gar-1.eps}


 }


 \end{center}


 \vspace*{-3pt}


     \Caption{Типовая структура СТО ИТС}


      \end{figure}


      


  Подсистема О$_1$ обеспечивает управление системой, производит сбор 


информации о состоянии элементов ИТС, формирует план распределения 


средств технического обеспечения и доводит вытекающие из него задачи до 


подсистем О$_2$ и О$_3$.


  


  Подсистема О$_2$ включает группировку сил и средств, привлекаемых к 


восстановлению работоспособности телекоммуникационной системы (ТКС).


  


  Подсистема О$_3$ выполняет функции эвакуации ТС, 


требующих сложных видов ремонта, обеспечение ЗИП и исправными 


ТС.


  


  Система~$E$ создает ограничения на эффективность функционирования СТО. 


  





  


  Приведенная на рис.~1 структурная схема СТО является в общем случае 


типовой для любого элемента подсистемы СТО, а поэтому правомерно 


многоуровневое представление модели ее функционирования.


  


  На рис.~2 приведена схема функционирования СТО ИТС, построенная по 


принципу иерархии, т.\,е.\ отражающая соподчиненность или согласованный по 


подчиненности порядок подсистем.


  


  Рассмотрим вопрос о моделировании функционирования СТО ИТС. При 


этом будем исходить из того положения, что функционирование СТО в первую 


очередь связано с распределением ресурсов на техническое обеспечение ее 


элементов.


  


  В общем виде выработка решения о распределении ресурсов СТО может 


быть представлена отображением


  \begin{equation}


  Q_1: R\times P\times S\times T\times Z\rightarrow V\,,


  \label{e1-gr}


  \end{equation}


где $r\hm\in R$~--- множество ресурсов СТО;


  $p\hm\in P$~--- множество объектов обеспечения;


  $s\hm\in S$~--- множество характеристик УЭ СТО ИТС;


  $t\hm\in T$~--- множество моментов времени;


  $z\hm\in Z$~--- множество целей, которые должны быть достигнуты в ходе 


технического обеспечения;


  $v\hm\in V$~--- множество решений о распределении ресурсов СТО по 


элементам ИТС.


  





  Отображение~(\ref{e1-gr}) представляет собой алгоритм, который каждому 


набору условий принятия решений ($r\hm\in R$, $p\hm\in P$, $s\hm\in S$, 


$z\hm\in Z$, $t\hm\in T$) ставит в соответствие некоторое решение $v\hm\in V$ 


из множества допустимых. Реализация этого решения приводит к некоторому 


результату, прогнозирование которого на этапе выработки решения может быть 


представлено отображением


  \begin{equation}


Q_2: R \times P \times S \times T \times V\rightarrow F\,,


\label{e2-gr}


\end{equation}


где $F$~--- множество возможных результатов технического обеспечения ИТС.


  


  Отображение~(\ref{e2-gr}) представляет собой модель прогнозирования 


результатов использования ресурсов СТО $r\hm\in R$ для технического 


обеспечения элементов ИТС $p\hm\in P$ в моменты времени $t\hm\in T$ в 


соответствии с решением $v\hm\in V$ в УЭ $s\hm\in S$.


  


  Качество решения определяется его конечным результатом при 


использовании имеющихся возможностей СТО.





 


  


  Оценка качества технического обеспечения может быть представлена 


отображением 


  \begin{equation*}


Q_3: R \times P \times S \times T\times F \rightarrow E\,,


%\label{e3-gr}


\end{equation*}


где $E$~--- упорядоченное по степени предпочтения множество оценок 


достижения поставленных целей $z\hm\in Z$. 





 \begin{figure} %fig2


\vspace*{1pt}


\begin{center}


 \mbox{%


 \epsfxsize=106.233mm


 \epsfbox{gar-2.eps}


 }


 \end{center}


 \vspace*{-10pt}


\Caption{Многоуровневое представление модели функционирования СТО ИТС}


\vspace*{-9pt}


\end{figure}





 Решение ищется при условии, что цели, включенные в 


отображение~(\ref{e1-gr}), являются измеримыми. 


  


  В общем случае СТО ИТС является управляемой системой, а поэтому 


дополним приведенную выше модель~(\ref{e1-gr}) моделью процесса 


управления. Воспользуемся для этого результатами работы~\cite{24-gr}.


  


  Примем для описания в качестве базовой иерархическую структуру с двумя 


уровнями управления. Особенностью такой структуры является сочетание 


централизованного управления системой с локальным управлением 


подсистемами с разделением функций и алгоритма управления, что полностью 


соответствует рассматриваемому случаю.


  


  Введем следующие обозначения:\\[-9pt]


  


  $m\in M$~--- множество управляющих сигналов;\\[-9pt]


  


  $\omega \in \Omega$~--- множество внешних воздействий;\\[-9pt]


  


  $l\in L$~--- множество информационных сигналов 1-го уровня; \\[-9pt]


  


  $k\in K$~--- множество информационных сигналов 2-го уровня;\\[-9pt]


  


  $u\in U$~--- множество координирующих сигналов;\\[-9pt]


  


  $y\in Y$~--- множество выходных сигналов.\\[-9pt]


  


  Тогда управляемый процесс $G$ на языке теории множеств можно задать 


отображением 


  \begin{equation*}


     G: M \times \Omega \rightarrow Y\,,


%     \label{e4-gr}


     \end{equation*}


где $M = M_1\times M_2\times\cdots \times M_i\times\cdots\times M_n$, а 


$M_i$~--- множество управляющих сигналов для $i$-го управляющего органа 


[$m_i\hm\in M_i$,  $m\hm=(m_1, m_2, \ldots , m_n)$].


  


  Модель функционирования $i$-й подсистемы управления ($i\hm = 1, 2, \ldots 


, n$) реализуется в виде отображения


  \begin{equation*}


     C_i: U \times L_i \rightarrow  M_i\,;\ L= L_1 \times L_2 \times \cdots \times  L_i 


\times \cdots \times L_n\,.


%     \label{e5-gr}


     \end{equation*}


     


     Модель службы управления системой реализуется в виде отображения


     \begin{equation*}


     C_0: K \rightarrow U\,.


%     \label{e6-gr}


     \end{equation*}


     


     Соответственно, информационные обратные связи 1-го и 2-го уровней 


управ\-ле\-ния реализуются отображениями


     \begin{align*}


     f_i: &\ \ \ M_i \times \Omega \times Y \rightarrow  L_i\,,\enskip  i = 1, 2, \ldots , 


n\,; %\label{e7-gr}


\\


     f_0: &\ \ \ U \times L \times M \rightarrow  K\,.  %\label{e8-gr}


     \end{align*}


  


  Дадим формализованное описание УЭ СТО ИТС.


  


  Отметим, что эксплуатация СТО происходит под влиянием случайных 


факторов, включающих воздействие пониженных и повышенных температур, 


вибрационные и ударные воздействия, пыль, а также возможно наличие 


неблагоприятных внешних воздействий (НВ): импульсного электромагнитного 


излучения, сейсмического удара, стихийных сил природы и других.


\pagebreak


  


  Введем следующие обозначения: 


  \begin{itemize}


\item   $s_{\mathrm{ВФ}}\in S_{\mathrm{ВФ}}$~--- множество внешних воздействующих 


факторов на элементы ИТС;


  \item


  $s_{\mathrm{НВ}}\in S_{\mathrm{НВ}}$~--- множество неблагоприятных 


воздействий.


\end{itemize}


  


  


  Тогда $S = S_{\mathrm{ВФ}}\cup S_{\mathrm{НВ}}$~--- условия эксплуатации СТО 


ИТС.


  


  Произведение $S=s^1_{\mathrm{ВФ}}\times s^2_{\mathrm{ВФ}} 


\times\cdots\times s^n_{\mathrm{ВФ}}\times s^1_{\mathrm{НВ}}\times 


s^2_{\mathrm{ НВ }} \times\cdots\times s^m_{\mathrm{ НВ }}$ будем называть 


общим пространством воздействий на элементы ИТС, на устойчивость к 


которым должна разрабатываться аппаратура, приборы и оборудование ИТС.





\begin{figure} %fig3


\vspace*{1pt}


\begin{center}


 \mbox{%


 \epsfxsize=127.463mm


 \epsfbox{gar-3.eps}


 }


 \end{center}


 \vspace*{-6pt}


\Caption{Общий методический подход к формированию СТО ИТС}


\end{figure}


  


  На рис.~3 представлен общий методический подход к формированию СТО 


ИТС, который может рассматриваться как дополнение к работам, изложенным 


выше.











  Для обеспечения функционирования СТО ИТС должна быть разработана 


система нормативных документов, регламентирующих: 


  \begin{itemize}


  \item цели, задачи и содержание работ, выполняемых в рамках СТО ИТС 


(под\-сис\-темы);


  \item требования к времени восстановления, сроку службы, сроку хранения и 


условиям эксплуатации ТС ИТС и СТО;


  \item состав и оснащенность СТО ремонтными мастерскими; 


  \item структуру и состав типовых комплектов ЗИП.


  \end{itemize}


  


  При этом статус документов, включающих изложенные выше вопросы, 


может быть различный, например: 


  \begin{itemize}


  \item  <<Положение о техническом обеспечении эксплуатации 


информационно-те\-ле\-ком\-муникационных сетей>>;


  \item  <<Регламент проведения ремонта и обслуживания технических 


средств>>; 


  \item  <<Положение об авторском надзоре над эффективностью 


функционирования системы технического обеспечения>>;


  \item  <<Программа оснащения системы технического обеспечения 


ремонтными мас\-тер\-скими>> и~т.\,д.


  \end{itemize}


  


\section{Заключение}





  В работе излагается последовательность работ по созданию СТО ИТС с 


учетом этапов жизненного цикла продукции. Рассмотрены основные вопросы, 


которые должны быть проработаны в ходе исследований, проектирования и 


эксплуатации. Выполнено моделирование функционирования СТО, дан вариант 


структуры СТО, приведено многоуровневое представление и общий 


методический подход к формированию СТО ИТС. Даны рекомендации по 


составу нормативного обеспечения функционирования СТО. 





 {\small\frenchspacing


{%\baselineskip=10.8pt


\addcontentsline{toc}{section}{Литература}


\begin{thebibliography}{99}





\bibitem{3-gr} %1


\Au{Дудник Б.\,Я., Овчаренко В.\,Ф., Орлов В.\,К. и~др.} Надежность и живучесть систем связи.~--- М.: 


Радио и связь, 1984. 216~с.





\bibitem{4-gr} %2


\Au{Дружинин Г.\,В.} Об оценке надежности сетей связи~/\!\!/ Надежность и контроль 


качества~/ Под ред. Б.\,В.~Гнеденко.~--- М.: Изд-во стандартов, 1984. Вып.~9. С.~10--16.





\bibitem{1-gr} %3


\Au{Филин Б.\,П.} Методы анализа структурной надежности сетей связи.~--- М.: Радио и 


связь, 1988. 204~с.


\bibitem{2-gr} %4


Надежность и живучесть радиоэлектронных систем.~--- Л.: ЦНИИ <<Румб>>, 1990. 123~с.





\bibitem{6-gr} %5


\Au{Ицкович А.\,А.} Построение системы управления эффективностью процесса 


технической эксплуатации самолетов~/\!\!/ Вопросы обеспечения эффективности 


технической эксплуатации летательных аппаратов.~--- М.: МИИГА, 1984. С.~9--19.


\bibitem{5-gr} %6


\Au{Ицкович А.\,А.} Управление эффективностью процесса технической эксплуатации 


машин.~--- М.: Знание, 1986. 106~с.





\bibitem{7-gr}


\Au{Буроменский Н.\,Г., Минцкер В.\,М.} Методологические вопросы  


эффективности мероприятий технического обеспечения войск~/\!\!/ Военная мысль, 


1993. №\,8. С.~49--55.


\bibitem{8-gr}


ГОСТ РВ 15.004-2004. Система разработки и постановки продукции на производство. 


Военная техника. Стадии жизненного цикла изделий и материалов.~--- М.: 


Госстандарт России, 2004. 23~с.








\bibitem{10-gr} %9


\Au{Буроменский Н.\,Г.} 


РД В 319.015-2006. Система добровольной сертификации РЭА, ЭРИ и материалов 


военного назначения <<Военэлектронсерт>>. Изделия электронной техники, 


квантовой электроники и электротехнические военного назначения. Требования к 


системе менеджмента качества.~--- Мытищи: 22 ЦНИИИ МО РФ, 2006.





\bibitem{11-gr} %10


ГОСТ ISO 9000-2011. Системы менеджмента качества. Основные положения и 


словарь.~--- М.: Стандартинформ, 2012. 32~с.





\bibitem{9-gr} %11


ГОСТ РВ 0015-002-2012. Система разработки и постановки продукции на производство. 


Военная техника. Системы менеджмента качества. Общие требования.~--- М.: 


Стандартинформ, 2012. 42~с.





\bibitem{12-gr} %12


\Au{Строгалев В.\,П., Толкачева И.\,О., Чуев В.\,Ю.}


Внешнее проектирование технических систем.~--- М.: МГТУ им.\ Н.\,Э.~Баумана, 1995. 72~с.


\bibitem{13-gr}


Руководство по техническому обеспечению связи и автоматизированных средств 


управления.~--- Мытищи: 16 ЦНИИИ МО РФ, 1987. 11~с.


\bibitem{14-gr}


\Au{Буроменский Н.\,Г., Минцкер В.\,М.} 


Ремонт радиоэлектронного вооружения в армии США~/\!\!/ Техника и вооружение, 1991. 


№\,6. С.~38--39.


\bibitem{15-gr}


\Au{Буроменский Н.\,Г., Кудашев Г.\,Н.} Критерии предельного состояния при 


определении долговечности военной техники связи.~--- Мытищи: 16 ЦНИИИ МО РФ, 


НТС №\,3, 1982.


\bibitem{16-gr}


\Au{Зацаринный А.\,А., Гаранин А.\,И., Козлов~С.\,В., Кондрашев~В.\,А.} Особенности 


расчета комплектов ЗИП в автоматизированных информационных системах в 


защищенном исполнении~/\!\!/ Системы и средства информатики, 2013. Т.~23. №\,1. 


С.~113--131.


\bibitem{17-gr}


\Au{Буроменский Н.\,Г.} Методические основы построения модели <<поражающее 


действие\,--\,стойкость>>~/\!\!/ Вопросы атомной науки и техники. Серия: Физика 


радиационного воздействия на радиоэлектронную аппаратуру, 2012. Вып.~4. 


С.~53--57.


\bibitem{18-gr}


\Au{Буроменский Н.\,Г.} Войска связи свою задачу выполнили~/\!\!/ Москва~--- 


Чернобылю. Кн.~1.~--- М.: Воениздат, 1998. 


С.~290--304.


\bibitem{19-gr}


ГОСТ РВ 15.203-2001. Система разработки и постановки продукции на производство. 


Военная техника. Порядок выполнения опытно-конструкторских работ по созданию 


изделий и их составных частей. Основные положения.~--- М.: Госстандарт России, 


2001. 117~с.


\bibitem{20-gr}


ГОСТ РВ 15.208-2005. Система разработки и постановки продукции на производство. 


Военная техника. Единый сквозной план создания образца (системы, комплекса) и его 


(их) составных частей. Общие положения.~--- М.: Стандартинформ, 2005. 21~с.





\bibitem{22-gr} %21


\Au{Зацаринный А.\,А., Буроменский Н.\,Г.} Сертификация систем и средств связи 


военного назначения~/\!\!/ Информационно-технический бюллетень НС ВС РФ, 1995. 


№\,14. С.~8--12.





\bibitem{21-gr} %22


РД В 319.35.04.00-2001. Положение о порядке применения электронных модулей, 


комплектующих изделий, электрорадиоизделий и конструктивных материалов 


иностранного производства в системах, комплексах и образцах вооружения и военной 


техники и их составных частях.~--- Мытищи: 22 ЦНИИ МО РФ, 2001.





\bibitem{23-gr}


\Au{Зацаринный А.\,А., Гаранин А.\,И., Ионенков Ю.\,С.} 


Методический подход к обоснованию требований к надежности 


информационно-телекоммуникационных сетей~/\!\!/ Системы и средства информатики.~--- 


М.: ИПИ РАН, 2010. Вып.~20. №\,3. С.~157--173.


\bibitem{24-gr}


\Au{Денисов А.\,А., Колесников Д.\,Н.} Теория больших систем управления.~--- 


Л.: Энергоиздат, 1982. 288~с.





 \end{thebibliography}


}


}





\vspace*{12pt}





\hrule





\vspace*{2pt}





\hrule





\def\tit{METHODOLOGICAL QUESTIONS OF~FORMATION OF~THE~SYSTEM TECHNICAL 


SUPPORT OF~INFORMATION-TELECOMMUNICATION NETWORKS}





\def\aut{A.\,A.~Zatsarinnyy$^1$, N.\,G.~Buromenskii$^2$, and A.\,I.~Garanin$^3$}








\titelr{\tit}{\aut}





\vspace*{3pt}





\noindent


$^1$IPI RAN, Moscow, Russia, azatsarinny@ipiran.ru\\


\noindent


$^2$Branch of the Federal State Government Institution ``46 Central


Reserach Institute\\


\hphantom{$^2$}of Defence Ministry of Russia,'' Mytishchi, Russia, buromienskii@mail.ru\\


\noindent


$^3$IPI RAN, Moscow, Russia, algaranin@mail.ru





%\vspace*{-12pt}


  


\Abste{The basic concepts in relation to the object of research are described. It is 


shown that the formation of the system of technical support (STS) of information


and 


telecommunication network (ITN) is an integrated system problem, acquiring\linebreak\vspace*{-12pt}


}





\Abstend{preconceptual covering of a very wide range of interrelated processes at all stages of the 


life cycle of ITN. Basic of them are formulated almost in the same sequence as they should be 


resolved in the course of creation of the STS. Consistently, the stages of 


research to create an STS, the tasks, and obtained results are described. The 


analytical model of the functioning of the STS and


general scheme of the method of 


forming the STS are given. It is shown that


the mathematical description of the processes of 


functioning of complex


hierarchical systems, such as STS of ITN, is extremely 


complex and serious problem. Its solution is suggested in two stages: at the first stage, 


the description is given in the form of abstract mathematical formalization of the 


process of functioning with the use of set theory (relations on the set between sets and 


functional maps), and at the second stage, the most important processes in the 


system, such as repair, maintenance, supply of new technical means, spare parts tools and accessories are


distinguished and 


transition from intuitive concepts to detailed mathematical description of using 


analytical and probabilistic models is realized. An approximate structural scheme of the 


STS is suggested. Basic content of the regulatory support of the STS of ITN is justified.}





  


\KWE{system; information and telecommunications system; the system of 


technical support; research; design; development; operation; effectiveness; modeling; 


structure}





\DOI{10.14357\!/08696527130211}





\renewcommand{\bibname}{\protect\rmfamily References}





 {\small\frenchspacing


{%\baselineskip=10.8pt


%\addcontentsline{toc}{section}{Литература}


\begin{thebibliography}{99}





\bibitem{3-gr-1} %1


\Aue{Dudnik, B.\,Ya., V.\,F.~Ovcharenko, V.\,K.~Orlov, \textit{et al.}} 1984. \textit{Nadezhnost' i zhivuchest' 


sistem svyazi} [\textit{Reliability and survivability of


communication systems}]. M.: Radio and Communications. 216~p.


\bibitem{4-gr-1} %2


\Aue{Druzhinin, G.\,V.} 1984. Ob ocenke nadezhnosti setej svyazi 


[On the assessment of reliability of  communication networks].


\textit{Nadezhnost' i  kontrol' kachestva}


[\textit{Reliability and quality control}]. Ed. B.\,V.~Gnedenko. 


M.: Publishing House of Standards. 9:10--16. 





\bibitem{1-gr-1} %3


\Aue{Filin, B.\,P.} 1988. \textit{Metody analiza strukturnoj nadezhnosti setej svyazi} 


[\textit{Methods of analysis of structural reliability of communication networks}]. M.: 


Radio and Communications. 204~p.


\bibitem{2-gr-1} %4


\textit{Nadezhnost' i zhivuchest' radiojelektronnyh sistem}. 


[\textit{Reliability and survivability of electronic systems}].  1990. Leningrad: 


Central Research Institute ``RUMB.'' 123~p.





\bibitem{6-gr-1} %5


\Aue{Ickovich, A.\,A.} 1984. Postroenie sistemy upravleniya jeffektivnost'yu processa 


tex\-ni\-ches\-koj jekspluatacii samoletov 


[Building efficiency management system of 


technical exploitation of aircrafts]. \textit{Voprosy obespecheniya jeffektivnosti 


texnicheskoj jekspluatacii letatel'nyx apparatov}  


[\textit{The issues of ensuring the effectiveness of the 


technical exploitation of aircrafts}]. M.: MIIGA [MIECA]. 9--19.





\bibitem{5-gr-1} %6


\Au{Ickovich, A.\,A.} 1986. \textit{Upravlenie jeffektivnost'yu processa texnicheskoj 


jekspluatacii mashin} [\textit{Performance management process of technical operation of 


vehicles}]. M.: Znanie.  106~p.





\bibitem{7-gr-1}


\Aue{Buromenskii, N.\,G., and V.\,M.~Mincker}. 1993. {Metodologicheskie voprosy 


ocenki jeffektivnosti meropriyatij texnicheskogo obespecheniya vojsk} 


[{Methodological issues of assessing the effectiveness of activities of the technical 


support for the troops}] \textit{Voennaya Mysl'} [\textit{Military Thought}] 8:49--55.


\bibitem{8-gr-1}


GOST RV 15.004-2004. 2004. Sistema 


razrabotki i postanovki produkcii na proizvodstvo. Voennaya texnika. Stadii zhiznennogo 


cikla izdelij i materialov. [Systems of development and launch of new products.


Military equipment. Stages of the life cycle of products and materials]. M.: 


Russian State Standards Publs. 23~p.








\bibitem{10-gr-1} %9


\Aue{Buromenskii, N.\,G.} 2006. RD V 319.015-2006. Sistema 


dobrovol'noj sertifikacii RJEA, JERI i materialov voennogo naznacheniya 


``Voenjelektronsert.'' Izdeliya jelektronnoj texniki, kvantovoj jelektroniki i 


jelektrotexnicheskie voennogo naznacheniya. Trebovaniya k sisteme menedzhmenta 


kachestva. [System of voluntary certification of REE, ERC, and materials for military


puposes ``Voenelectroncert'']. Mytishchi: 22 Central Research


Institute of Defence Ministry of Russia.





\bibitem{11-gr-1} %10


GOST ISO 9000-2011. 2012. Sistemy 


menedzhmenta kachestva. Osnovny'e polozheniya i slovar'. [Quality management systems. Fundamentals


and vocabulary]. 


M.: Standardinform Publs. 32~p.





\bibitem{9-gr-1} %11


GOST RV 0015-002-2012. 2012. Sistema 


razrabotki i postanovki produkcii na proizvodstvo. Voennaya texnika. Sistemy 


menedzhmenta kachestva. Obshhie trebovaniya.  [Systems of development and launch of new


 products. Military equipment. Quality management systems. General requirements].


 M.: Standardinform Publs. 42~p.


 





\bibitem{12-gr-1}


\Aue{Strogalev, V.\,P., I.\,O.~Tolkacheva, and V.\,Yu.~Chuev}. 1995. \textit{Vneshnee 


proektirovanie texnicheskix system}  


[\textit{The external design of technical systems}]. M.: 


Publishing House of the Moscow State University named 


N.\,E.~Bauman. 72~p.





\bibitem{13-gr-1}


16 CNIII 


MO RF [16~Central Research Institute of Defense Ministry of Russia].


1987. Rukovodstvo po texnicheskomu obespecheniyu svyazi i avtomatizirovannyx sredstv 


upravleniya [Guide to technical communication and automated management tools]. 


Mytishchi. 11~p.





\bibitem{14-gr-1}


\Aue{Buromenskii, N.\,G., and V.\,M.~Mincker}. 1991. Remont radiojelektronnogo vooruzheniya 


v armii SShA [Repair of electronic weapons in the U.S. army]. \textit{Texnika i 


Vooruzhenie} [\textit{Equipment and Arms}]  6:38--39. 





\bibitem{15-gr-1}


\Aue{Buromenskii, N.\,G., and G.\,N.~Kudashev}. 1982. Kriterii predel'nogo sostoyaniya pri 


opredelenii dolgovechnosti voennoj texniki svyazi [Limiting state criteria in determining the 


durability of military communications]. Scientific and technical collection ser.


Mytishchi: 16 Central Research Institute of 


Defense Ministry of Russia. 3.





\bibitem{16-gr-1}


\Aue{Zatsarinny, A.\,A., A.\,I.~Garanin, S.\,V.~Kozlov, and V.\,A.~Kondrashev}. 2013. 


Oso\-ben\-nosti rassheta komplektov ZIP v avtomatizirovannyx informacionnyx sistemax v 


zashhishhennom ispolnenii [About some particularities of spare parts sets calculation for secured


information systems]. 


\textit{Systems and Means of Informatics}  23(1):113--31.


\bibitem{17-gr-1}


\Aue{Buromenskii, N.\,G.} 2012. Metodicheskie osnovy postroeniya modeli ``porazhayushhee 


dejstvie\,--\,stojkost'\,'' [Methodological bases of construction of model ``striking 


action\,--\,resistance'']. \textit{Voprosy atomnoj nauki i texniki}. 


Fizika radiacionnogo vozdeistviya na radiojelectronnuyu apparaturu ser.


[\textit{Issues of atomic science 


and technology}. Phisics of radiation exposure to radioelectronic equipment ser.]


4:53--57.





\pagebreak





\bibitem{18-gr-1}


\Aue{Buromenskii, N.\,G.} 1998. Vojska svyazi svoyu zadachu vypolnili 


[Signal corps completed their task]. 


\textit{Moskva~--- Chernobylyu}. [\textit{Moscow to Chernobyl}].


M.: Voenizdat. I:290--304.


\bibitem{19-gr-1}


GOST RV 15.203-2001. 2001. Sistema 


razrabotki i postanovki produkcii na proizvodstvo. Voennaya texnika. Poryadok vypolneniya 


opytno-konstruktorskix rabot po sozdaniyu izdelij i ix sostavnyx chastej. Osnovnye 


polozheniya [State Military Standard of Russian Federation. How to perform development work to


create products and their components. Fundamentals].  


M.: Russian State Standards Publs. 117~p.





\bibitem{20-gr-1}


GOST RV 15.208-2005. 2005. Sistema 


razrabotki i postanovki produkcii na proizvodstvo. Voennaya texnika. Edinyj skvoznoj plan 


sozdaniya obrazca (sistemy, kompleksa) i ego (ix) sostavnyx chastej. Obshhie polozheniya


[System of development and launch of new products. Military equipment.


Through and single plan to create a sample (system, complex) and their


components. General provisions].


M.: Standardinform Publs. 21~p.





\bibitem{22-gr-1} %21


\Aue{Zatsarinnyy, A.\,A., and N.\,G.~Buromenskii}. 1995. Sertifikaciya sistem i sredstv svyazi 


voennogo naznacheniya [Certification of communication means and systems for military 


use]. \textit{Informacionno-texnicheskij byulleten' NS VS RF} [\textit{Information and 


Technical Bulletin of the Russian Federation}]  14:8--12.





\bibitem{21-gr-1} %22


RD V 319.35.04.00-2001. 2001. Polozhenie o poryadke 


primeneniya jelektronnyx modulej, komplektuyushhix izdelij, jelektroradioizdelij i 


konstruktivnyx materialov inostrannogo proizvodstva v sistemax, kompleksax i obrazcax 


vooruzheniya i voennoj texniki i ix sostavnyx chastyax


[Thesis  on the application of electronic modules, components, electrical and radio devices and


structural materials of foreign production systems in systems and models of


weapons and military equipment and their components.] Mytishchi:


22~Central Research  Institute of Defense Ministry of Russia.





\bibitem{23-gr-1}


\Aue{Zatsarinnyy, A.\,A., A.\,I.~Garanin, and Yu.\,S.~Ionenkov}. 2010. Metodicheskij podxod k obosnovaniyu 


trebovanij po nadyozhnosti k sostavnym chastyam informacionno-telekommunikacionnyx 


setej [Methodical approach to substantiation of requirements to reliability to the parts of 


information-telecommunication networks]. 


\textit{Systems and Means of Informatics} 20(3):157--73.








 \label{end\stat}


 


 \bibitem{24-gr-1}


\Aue{Denisov, A.\,A., and D.\,N.~Kolesnikov}. 1982. 


\textit{Teoriya bol'shix sistem upravleniya} 


[\textit{Theory of large systems management}].


Leningrad: Energoizdat, Leningrad  Department.  288~p.


  


 \end{thebibliography}


} }





\renewcommand{\bibname}{\protect\rmfamily Литература}


  


